\documentclass[10pt,a4paper]{amsart}
\usepackage[margin=19mm]{geometry}
\usepackage{amsmath,amssymb,mathtools}
\usepackage[scr=rsfs,cal=euler]{mathalfa}
\usepackage{xfrac}
\usepackage[unicode,hidelinks,bookmarks=false]{hyperref}
\newcommand{\ie}{i.e.\ }
\newcommand{\cf}{cf.\ }
\newcommand{\resp}{respectively\ }
\newcommand{\Subsection}{\subsection}
\providecommand{\nobreakdash}{\nobreak\hbox{-}}
\setlength{\emergencystretch}{2em}
\multlinegap0pt
\usepackage{bm}
\usepackage{bbm}
\usepackage{dsfont}
\usepackage{xspace}
\usepackage{cancel}
\usepackage{mathtools}
\usepackage{amsthm}
\makeatletter
\@ifundefined{theorem}{\newtheorem{theorem}{Theorem}}{}
\@ifundefined{corollary}{\newtheorem{corollary}[theorem]{Corollary}}{}
\makeatother
\newcommand{\ZZ}{{\mathbb Z}}
\title[Cutting a convex body into fat parts and approximating Eucli]{Cutting a convex body into fat parts and approximating Euclidean distance by graph distances}
\author{J\'anos Pach, G\'abor Tardos}
\date{Source: arXiv:2609.20702v1 (CC BY 4.0)}
\begin{document}
\maketitle

\noindent\emph{Source and licence.} Cutting a convex body into fat parts and approximating Euclidean distance by graph distances. Version arXiv:2609.20702v1 (CC BY 4.0), distributed under the Creative Commons Attribution 4.0 International licence, as recorded in the arXiv licence metadata for this exact version and shown on the abstract page https://arxiv.org/abs/2609.20702v1. Full rights notice: licenses/arxiv-2609.20702-CC-BY-4.0.txt.\par\smallskip

\noindent\textit{Editorial scope.} Counted unit: the Corollary whose source label is \texttt{inwheelongrid} in the article (source0.tex lines 125--128, source label \texttt{inwheelongrid}), delivered together with the article's own complete original proof of it (source0.tex lines 130--168, one \texttt{proof} environment of 39 lines). No mathematics is written, repaired or re-derived: statement and proof are the article's own text line for line. Required context retained uncounted: none. Every definition, display and label of those lines is retained unchanged. Omitted from this package and not redistributed: 1--124, 129--129, 169--728. Nothing inside a retained excerpt is omitted or altered: every retained body is delivered exactly as the article's lines are written, so no omission receipt of any kind accompanies this package. Citations: the retained excerpts carry no citation command at all, so no bibliography entry of the article is transcribed and no reference is redistributed. Reference adaptation: none was needed and none was made; every reference inside the delivered unit resolves inside it, so no unresolved reference or citation is produced. Printed slips kept literal: no printed word, symbol, hypothesis or number of the counted statement or of its proof was altered. Layout and class: the article's own class is \texttt{article} with packages including amssymb, amsmath, amsthm, graphicx, comment, geometry, todonotes, caption, subcaption, overpic, pgfplots, mathrsfs, tikz; the wrapper is the lane's pinned \texttt{amsart} editorial wrapper at 10pt on A4, which restates the theorem-like environments the retained text needs and adds the article's own macro definitions that the retained text needs (\texttt{\textbackslash ZZ}), with extra wrapper packages (bm, bbm, dsfont, xspace, cancel, mathtools). The delivered numbering is the unit's own. Assets: no retained span contains a figure environment, a graphics-inclusion command, a PDF-page inclusion, a table or a TikZ picture; no image file is copied into this package, the package declares no asset dependency and no image file is redistributed with it. Licence: version arXiv:2609.20702v1 of this article is distributed under the Creative Commons Attribution 4.0 International licence, as recorded in the arXiv licence metadata for this exact version and shown on the abstract page https://arxiv.org/abs/2609.20702v1. A full rights notice is delivered at licenses/arxiv-2609.20702-CC-BY-4.0.txt. Characters: every retained excerpt is pure ASCII, so the delivered engine, XeLaTeX with its default Latin Modern text font, reports no missing character.
\begin{corollary}\label{inwheelongrid}
There exists a graph $H'$ on the vertex set ${\ZZ}^2$ such that for any two integer points $p$ and $q$, we have $$d_{H'}(p,q)=d(p,q)+O\left(\frac{d(p,q)}{\log^{2/3}(d(p,q))}\right),$$
as $d(p,q)$ tends to infinity.
\end{corollary}

\begin{proof}[Proof of Corollary~\ref{inwheelongrid}]
To construct $H'$, we take a point set $P$ described in the theorem, blow it up by a sufficiently large constant factor, and then embed it in the grid. For a precise argument, we first bound the maximum degree of the graph $H$ from the theorem. The neighbors of a vertex $x\in P$ in the graph $H$ are all within distance $1$ from $x$, but are at distance at least $1/3$ from each other. Therefore, the open disks of radius $1/6$ centered at the neighbors of $x$ are pairwise disjoint and are contained in the disk of radius $7/6$ centered at $x$. Comparing areas, we obtain
$$
\deg_H(x)<\frac{\pi(7/6)^2}{\pi(1/6)^2}=49.
$$
Thus, the maximum degree of $H$ is at most $48$.

To obtain $H''$, replace each edge of $H$ by a path of length $501$. The new graph $H''$ contains the vertices in $P$ and also new degree $2$ vertices on the connecting paths. Every vertex $v$ of $H''$ is within $H''$-distance at most $250$ from a unique vertex $x\in P$.

Now we injectively map the vertices $v$ of $H''$ to points $f(v)$ of the grid $\ZZ^2$ in such a way that, if $x\in P$ is the unique vertex satisfying
$d_{H''}(v,x)\leq 250,$ then $d(f(v),501x)<501/6.$ This is possible because

(i) every vertex $v$ of $H''$ is associated with exactly one vertex $x\in P$;

(ii) there are at most $48\cdot250+1=12001$
vertices $v$ associated with any fixed point $x\in P$;

(iii) the open disks of radius $501/6$ around the points in $501P$ are pairwise disjoint, since distinct points of $P$ are at distance at least $1/3$ from each other; and

(iv) every open disk of radius $501/6$ contains more than $12001$ grid points. Indeed, such a disk contains an open axis-parallel square of side length $118$, because $59\sqrt{2}<501/6.$

Every open interval of length $118$ contains at least $117$ integers, so this square contains at least
$117^2=13689>12001$ points of $\ZZ^2$.

Finally, we construct $H'$ on the vertex set $\ZZ^2$ by connecting $f(v)$ and $f(w)$ for every edge $vw$ of $H''$, and adding a single additional edge for every grid point $z$ \emph{not} of the form $f(v)$. For such a point $z\in\ZZ^2$, we select $x\in P$ within distance $1/3$ of $z/501$ and add the edge $zf(x)$. The vertices of $\ZZ^2$ that are not of the form $f(v)$ have degree $1$, and therefore they cannot create shortcuts between vertices in $f(V(H''))$. It follows that, for any $x,y\in P$,
$$
d_{H'}(f(x),f(y))=501d_H(x,y).
$$
Furthermore,
$$
d(f(x),f(y))=501d(x,y)+O(1).
$$
For every $z\in\ZZ^2$, there exists $x\in P$ such that
$d(z,f(x))=O(1)$ and $d_{H'}(z,f(x))=O(1).$ Indeed, if $z$ is not of the form $f(v)$, this follows from the definition of the additional edge incident to $z$. If $z=f(v)$, take the unique $x\in P$ satisfying $d_{H''}(v,x)\leq250$.

Using part~3 of the theorem and increasing the implicit constant, if necessary, to cover bounded distances, we obtain, for $d(z,z')\geq 2$,
$$d_{H'}(z,z')=d(z,z')+O\left(\frac{d(z,z')}{\log^{2/3}(d(z,z'))}\right).$$
This proves the corollary.
\end{proof}

\end{document}
